\chapter{Estado da Arte e Análise de Requisitos}
\label{cap2}

%Mais umas citações \cite{AndroidDoc, Chen1976}

%a linha seguinte deve ser substituída pelo texto do capítulo
A qualidade do ar cada vez mais foi comprovada ser um assunto de extrema importancia, tendo variaveis que enriquecem ou desfavorecem a saúde humana, mais concretamente respiratoria. Para dar inicio a pesquisa sobre este topico procurei encontrar estudos que contivessem pontos que a meu criterio pudessem ser validos e importantes para compreender o problema como um todo. Encontrei os seguintes relatorios tecnicos e plataforma institucional de dados meteorológicos, que me auxiliaram a compreender aquilo que seria este problema. 
O relatório publicado pela Agência Europeia do Ambiente em 2024 \cite{EEA2024} evidenciou os efeitos da exposição prolongada à poluição atmosférica na saúde pública, apresentando estimativas de mortalidade e morbilidade associadas a poluentes como partículas finas (PM\textsubscript{2,5}), dióxido de azoto (NO\textsubscript{2}) e Ozono (O\textsubscript{3}) . Neste artigo consegui compreender como a poluição afeta a saude humana, ou seja, neste ponto compreendi a \textbf{fundamentação do problema}.

Por sua vez, o relatório de 2026 \cite{EEA2026} demonstrou que, apesar das melhorias registadas na qualidade do ar europeu ao longo das últimas décadas, persistem situações de exposição a concentrações superiores aos valores recomendados pela Organização Mundial da Saúde. Estes resultados reforçam a importância da monitorização ambiental e da disponibilização de informação que permita acompanhar a evolução dos poluentes. Este documento permitiu perceber o estado atual da qualidade do ar na Europa, pelo que representa a \textbf{contextualização e a necessidade de monitorização}.

Complementarmente, a consulta da plataforma de observação meteorológica do \cite{IPMA} permitiu identificar diferentes variáveis relevantes para a caracterização das condições atmosféricas, nomeadamente temperatura, humidade relativa, pressão atmosférica e vento, bem como observar métodos de apresentação gráfica da informação. Por sua vez neste artigo, compreendi qual é o estado atual europeu, em relação à qualidade do ar, por sua vez esta foi a contextualização e necessidade para a monitorização análise das variáveis e da sua visualização.

Em conjunto, estas fontes constituem uma base para a fundamentação do presente projeto, demonstrando a relevância da recolha, tratamento, armazenamento e visualização de dados ambientais. Neste contexto, o desenvolvimento de uma aplicação web para análise da qualidade do ar na região de Beja procura contribuir para uma interpretação mais acessível e organizada das concentrações de poluentes e das variáveis meteorológicas disponibilizadas pelos sensores. Por fim, mas não menos importante. compreendi neste ponto como são observador e apresentados os dados metereologicos. Ou seja, a analise de variaveis de visualização.

\section{Contextualização do Projeto}
\label{sec:contextualizacao}

% Apresentação da Bettair
% Origem dos dados
% Contextualização do problema em Beja
% Objetivo e contribuição da aplicação
A informação presente neste projeto tem origem na plataforma Bettair, aqui é onde estão presente os dados recolhidos pelos sensores de qualidade do ar. Cada ficheiro disponibilizado corresponde a um sensor e contém as concentrações dos poluentes medidos por esse sensor, associadas a um intervalo de tempo.

O objetivo da aplicação é permitir ao cliente consultar, de forma organizada, os dados dos sensores instalados na região de Beja, relacionando as concentrações de poluentes com as condições meteorológicas registadas no mesmo intervalo. A aplicação possibilita a consulta temporal, isto dentro do periodo dos dados inserios, por exemplo de junho a agosto, só é possivel, ver os dados presentes neste intervalo de tempo.

\section{Fatores Analisados}

Para compreender o problema foi necessário, antes de mais, compreender os dados. No primeiro contacto com o ficheiro CSV, observou-se a estrutura apresentada na Tabela~\ref{tab:campos}, composta por 19 campos: quatro marcas temporais, treze medições (poluentes e variáveis meteorológicas) e as coordenadas geográficas do sensor.


\begin{table}[htbp]
\centering
\caption{Campos do ficheiro CSV disponibilizado pela Bettair.}
\label{tab:campos}

\small
\renewcommand{\arraystretch}{1.2}

\begin{tabularx}{\textwidth}{
    @{}
    >{\raggedright\arraybackslash}X
    >{\raggedright\arraybackslash}p{0.24\textwidth}
    >{\raggedright\arraybackslash}p{0.16\textwidth}
    @{}
}

\toprule
\textbf{Campo} & \textbf{Unidade} & \textbf{Tipo} \\
\midrule

timestamp\_init\_UTC
    & ISO 8601 (UTC) & Tempo \\

timestamp\_end\_UTC
    & ISO 8601 (UTC) & Tempo \\

timestamp\_init\_Europe/Lisbon
    & ISO 8601 (Lisboa) & Tempo \\

timestamp\_end\_Europe/Lisbon
    & ISO 8601 (Lisboa) & Tempo \\

NO2
    & \si{\micro\gram\per\cubic\metre} & Poluente \\

O3
    & \si{\micro\gram\per\cubic\metre} & Poluente \\

CO
    & \si{\milli\gram\per\cubic\metre} & Poluente \\

SO2
    & \si{\micro\gram\per\cubic\metre} & Poluente \\

PM1, PM10, PM2P5
    & \si{\micro\gram\per\cubic\metre} & Partículas \\

pressure
    & \si{\micro\gram\per\cubic\metre}$^{*}$
    & Meteorologia \\

relativeHumidity
    & \% & Meteorologia \\

temperature
    & \si{\celsius} & Meteorologia \\

windDirection
    & graus & Meteorologia \\

windSpeed, windGust
    & \si{\metre\per\second} & Meteorologia \\

latitude, longitude
    & graus decimais & Localização \\

\bottomrule

\end{tabularx}

\vspace{0.2cm}

\begin{minipage}{\textwidth}
\footnotesize
$^{*}$ A unidade indicada no cabeçalho é incoerente
com valores como 984,5, típicos de hPa
(ver Secção~3.1.2).
\end{minipage}

\end{table}

Através desta análise foi possível identificar os campos que teriam de ser tratados para dar início ao processo ETL (\textit{Extract, Transform and Load}), nomeadamente: cabeçalhos com unidades incorporadas no nome, marcas temporais em dois fusos horários, valores em falta e a unidade duvidosa da pressão atmosférica.









